\scanpage{109}
\mspage{108}
On dispose donc de deux généralisations possibles, en \unsure{direction} d'une topologie sans axiome de divisibilité. Mais c'est la direction $\unsure{|\circ|\,|\,|o|}$ qui me semble la plus riche. En effet, elle a toujours une structure de $\unsure{\gtrless|\lessgtr}$-\unsure{magasin|magma} correspondante, mais \unsure{avec}, \uline{en plus}, une notion $|\circ|$ de disjonction intrinsèque, qui exprime une intuition géométrique \unsure{palpable}, laquelle échappe (en l'absence de At L1) au seul $\gtrless$-\unsure{point} de vue. L'axiome axiome At C (\unsure{spec|sup}), \unsure{avec} le \unsure{recul}, apparaît artificiel -- il exprime l'embarras de celui qui ne dispose d'une notion de ``disjonction'' en termes des \unsure{illusoires} ``strates ouvertes'' $X^{\unsure{o|0|\circ}}$.

Donc finalement, je change d'optique -- la ribambelle d'axiomes de tout poil At C (spec), At C, At D\ldots\ \unsure{tombe} \unsure{en} un \unsure{bruit} \unsure{floche|flou}, Je reprends
